\scanpage{132}
% archive foliation: 131
À un tel $L$, on associe la structure complexe
\inserted{\deleted{\unsure{conjuguée}} \deleted{\unsure{directe}}, \deleted{conjuguée de celle déduite de}}
\deleted{$V$ \unsure{directe} de celle de $L$ par transport de}
$V$ qui est conjuguée de celle déduite de la str. complexe de $L$ par transport de structure via (7), i.e. on pose
\[
\tag{8}
ux=\varphi(-i\varphi^{-1}(x))=-\varphi(i\varphi(x))
\]
\margin{$u$ étant la multiplication par $i$ pour la structure complexe envisagée sur $V$}
\margin{cette convention est imposée par les \unsure{structures} de ``\unsure{modules}'' de courbes elliptiques \unsure{et de VA)}, du fait qu'on veut que les param. modulaires \unsure{qu'on} \unsure{(*)} introduit soient \uline{holomorphes} et non antiholomorphes \ldots}

Inversement, si $V$ est muni d'une structure complexe, \deleted{l'applic} faisant de $V$ un espace vectoriel complexe $V_c$, alors l'application \inserted{identique} $\mathbb R$-linéaire $V\longrightarrow V_c$ se prolonge, par la propriété universelle de $V_{\mathbb C}$, en une application $\mathbb C$-linéaire
\[
\tag{9}
\psi:V_{\mathbb C}\longrightarrow V_c
\]
de sorte que
\[
\tag{10}
\left\{\begin{aligned}
\psi(ix)&=ux\\
\psi(x)&=x
\end{aligned}\right.
\qquad\forall x\in V=V_{\mathbb R},
\]
donc
\[
\tag{11}
\psi(x+iy)=x+uy.
\]
\deleted{Posons} \deleted{Considérons} Considérons
\[
\tag{12}
\begin{aligned}
L^{\unsure{\prime|!}}=\operatorname{Ker}\psi
&=\{x+iy\mid x,y\in V_{\mathbb R},\ x+uy=0\ \text{i.e. }y=\unsure{u}x\}\\
&=\{\deleted{\unsure{x+iu(x)}}\mid x\in V_{\mathbb R}\}
\end{aligned}
\]
et considérons l'application correspondante (6)
\[
\begin{gathered}
\varphi:\underset{\substack{\Vert\\x+iu(x)}}z\longmapsto\tfrac12(z+\bar z)
=\Re z=x\deleted{\unsure{+uy}}:\quad L\longrightarrow V .
\end{gathered}
\]
On voit alors que cette application est bijective, et \uuline{antilinéaire},
\inserted{et que la \unsure{correspondance} à la structure complexe}
\deleted{\unsure{permet de récupérer la structure complexe primitive sur $V_{\mathbb R}$}},
\deleted{$L=L^{\unsure{\prime|!}}=\operatorname{Ker}\psi\ \unsure{\subset V_{\mathbb C},\ V=L_{\mathbb R}}$},
\margin{\unsure{voir une autre explication}}

(13) $[$On peut dire aussi, que la donnée d'une structure complexe sur $V$ équivaut :

1°) La donnée d'un \deleted{\unsure{gr}} sous-groupe compact \deleted{maximal} $U$ \deleted{de} \deleted{$\unsure{\mathrm{Gl}}$} $\in\operatorname{Aut}_{\mathbb R}(V)$, \deleted{i.e. d'un}
